\documentclass[11pt,a4paper]{article}
\usepackage[utf8]{inputenc}
\usepackage[T1]{fontenc}
\usepackage{lmodern}
\usepackage[spanish]{babel}
\usepackage[margin=2cm]{geometry}
\usepackage{booktabs}
\usepackage{longtable}
\usepackage[hidelinks]{hyperref}
\title{\textbf{BOLETA DE VOTACI\'{O}N PONDERADA}}
\author{TA-IND-02 -- Cl\'{i}nica veterinaria\\
Estudiante: Jeanpierre Robinson Espinoza}
\date{}
\begin{document}
\maketitle

\noindent Votante: C\'{o}digo \rule{2.2cm}{0.4pt}\quad Fecha: \rule{2.2cm}{0.4pt}

\noindent Distribuya \textbf{exactamente 100 puntos} entre los veinte requisitos. Use n\'{u}meros enteros no negativos. No escriba nombres.

\vspace{0.3cm}

\begin{longtable}{@{}lp{13.5cm}r@{}}
\toprule
ID & Enunciado & Puntos \\
\midrule
\endhead
R01 & El sistema deberá permitir registrar cada mascota con un código único y los campos obligatorios nombre, especie y propietario; además podrá guardar raza, sexo, edad, peso y señas particulares. & \\
R02 & El sistema deberá permitir registrar los datos de contacto del propietario, al menos nombre y teléfono, asociados a una o varias mascotas. & \\
R03 & El sistema deberá permitir registrar antecedentes, alergias y enfermedades previas de cada mascota. & \\
R04 & El sistema deberá permitir registrar cada consulta con los seis campos definidos: fecha, motivo, síntomas, diagnóstico, tratamiento y observaciones del veterinario. & \\
R05 & El sistema deberá mostrar en orden cronológico todas las consultas anteriores de una mascota. & \\
R06 & El sistema deberá permitir registrar las vacunas aplicadas a cada mascota, indicando tipo de vacuna, fecha de aplicación y responsable. & \\
R07 & El sistema deberá permitir registrar la fecha de la próxima dosis o del siguiente control de vacunación. & \\
R08 & El sistema deberá mostrar al personal de la clínica un aviso visible al menos 3 días antes de la fecha programada de vacunación, cuando exista una próxima dosis registrada. & \\
R09 & El sistema deberá mostrar una lista de mascotas cuya fecha de vacunación sea igual o anterior a los 7 días siguientes, para planificar los controles de la semana. & \\
R10 & El sistema deberá permitir buscar mascotas por nombre, código de paciente o teléfono del propietario. & \\
R11 & El sistema deberá gestionar citas veterinarias registrando fecha, hora, mascota y motivo de atención en un único registro consultable. & \\
R12 & El sistema deberá registrar el seguimiento complementario de cada mascota (desparasitación y control de peso) como eventos fechados con observación. & \\
R13 & El sistema deberá restringir la creación, modificación y eliminación de historias clínicas según el rol del usuario. & \\
R14 & El sistema deberá garantizar la trazabilidad de las modificaciones de historias clínicas, registrando el usuario responsable y la fecha y hora de cada cambio. & \\
R15 & El sistema deberá mantener la integridad de los registros evitando duplicados por código de mascota o por coincidencia de nombre, especie y teléfono del propietario. & \\
R16 & El sistema deberá permitir consultar una ficha clínica en máximo 3 segundos en el 95 \% de los casos, con hasta 1.000 mascotas registradas y conexión estable. & \\
R17 & El sistema deberá realizar al menos una copia de seguridad diaria de la información clínica y permitir recuperar los datos respaldados en máximo 4 horas. & \\
R18 & El sistema deberá permitir que al menos 4 de cada 5 usuarios del personal completen el registro y la búsqueda de una mascota en máximo 3 minutos, sin asistencia, en una prueba de uso. & \\
R19 & El sistema deberá mantener la portabilidad del historial clínico, permitiendo exportarlo o imprimirlo con datos completos y legibles para entregarlo al propietario o compartirlo con otro profesional. & \\
R20 & El sistema deberá mantener disponibles los registros durante el horario de atención definido y configurable por la clínica, con al menos 99 \% mensual. & \\
\midrule
 & \textbf{Total} & \textbf{100} \\
\bottomrule
\end{longtable}

\noindent Comprobaci\'{o}n: la suma de los veinte valores debe ser 100. Si no lo es, corrija sus propios valores.
\end{document}

